\documentclass[10pt]{beamer}
\usetheme{metropolis}
\usepackage{graphicx}
\usepackage{listings}
\usepackage{booktabs}
\usepackage{xcolor}
\usepackage{hyperref}

\definecolor{codebg}{HTML}{F5F5F5}
\definecolor{codegreen}{HTML}{2E7D32}
\definecolor{codegray}{HTML}{757575}
\definecolor{codeblue}{HTML}{1565C0}

\lstdefinestyle{rstyle}{
  language=R, backgroundcolor=\color{codebg},
  basicstyle=\ttfamily\scriptsize,
  keywordstyle=\color{codeblue}\bfseries,
  stringstyle=\color{codegreen},
  commentstyle=\color{codegray}\itshape,
  breaklines=true, frame=single,
  rulecolor=\color{codegray!30}, showstringspaces=false, tabsize=2,
  morekeywords={library,ggplot,aes,geom_col,geom_text,geom_line,geom_point,
    coord_flip,theme_minimal,theme,element_text,element_blank,labs,
    fct_reorder,scale_fill_manual,facet_wrap,mutate}
}
\lstdefinestyle{pythonstyle}{
  language=Python, backgroundcolor=\color{codebg},
  basicstyle=\ttfamily\scriptsize,
  keywordstyle=\color{codeblue}\bfseries,
  stringstyle=\color{codegreen},
  commentstyle=\color{codegray}\itshape,
  breaklines=true, frame=single,
  rulecolor=\color{codegray!30}, showstringspaces=false, tabsize=4,
  morekeywords={import,from,as,pd,plt,sns,np,fig,ax}
}

\title{Critique \& Redesign Workshop}
\subtitle{Workshop 1 --- Module 7}
\author{Prof.Asoc.\ Endri Raco}
\institute{Data Visualization for Data Scientists \\ UNYT Tirana}
\date{2025/26}

\begin{document}

\begin{frame}
  \titlepage
\end{frame}

% ================================================================
\begin{frame}{Module Roadmap}
  \begin{enumerate}
    \item The four-quadrant critique framework
    \item Five canonical redesign case studies
    \item The redesign workflow: identify → diagnose → fix → validate
    \item Peer review protocol for visualization
  \end{enumerate}
  \begin{exampleblock}{Learning Outcome}
    You will be able to systematically critique any chart against four
    quality dimensions and execute a principled redesign using the
    techniques from Modules 1--6.
  \end{exampleblock}
\end{frame}

% ================================================================
\section{The Critique Framework}
% ================================================================

\begin{frame}{Four Quadrants of Critique}
  \begin{center}
    \includegraphics[width=0.82\textwidth]{figures_m07/critique_framework}
  \end{center}
\end{frame}

% ================================================================
\begin{frame}{Using the Framework}
  \centering
  \begin{tabular}{lp{4cm}l}
    \toprule
    \textbf{Quadrant} & \textbf{Key Question} & \textbf{Module Ref.} \\
    \midrule
    Data Integrity    & Does it honestly represent the data? & M01, M02 \\
    Visual Encoding   & Is each variable on the best channel? & M03, M04 \\
    Design Craft      & Is the typography/colour/whitespace clean? & M04, M05 \\
    Communication     & Does the audience get the message? & M05, M06 \\
    \bottomrule
  \end{tabular}

  \vspace{8pt}
  \begin{alertblock}{Pro-Tip}
    Critique the chart, not the author. Use the framework as an objective
    diagnostic tool: ``The lie factor is 3.5 because the axis starts at
    95'' is specific and constructive; ``this chart is bad'' is not.
  \end{alertblock}
\end{frame}

% ================================================================
\section{Five Redesign Case Studies}
% ================================================================

\begin{frame}{The Redesign Workflow}
  \begin{center}
    \includegraphics[width=0.92\textwidth]{figures_m07/redesign_workflow}
  \end{center}
  \vspace{4pt}
  Every redesign follows five steps: identify the problem, diagnose
  which quadrant it belongs to, sketch an alternative, implement in
  code, and validate with a peer.
\end{frame}

% ================================================================
\begin{frame}{Redesign 1: 3-D Pie $\rightarrow$ Horizontal Bar}
  \begin{center}
    \includegraphics[width=0.92\textwidth]{figures_m07/redesign_pie}
  \end{center}
  \begin{exampleblock}{Data Insight}
    \textbf{Diagnosis}: Encoding (angle) + Craft (chartjunk: shadow, explosion). \\
    \textbf{Fix}: Replace angle with position (horizontal bar), sort by value,
    add direct labels, remove all decoration.
  \end{exampleblock}
\end{frame}

% ================================================================
\begin{frame}[fragile]{Redesign 1 in R}
  \begin{columns}[T]
    \begin{column}{0.52\textwidth}
\begin{lstlisting}[style=rstyle]
library(tidyverse)

df <- tibble(
  dept = c("Marketing","Engineering",
    "Sales","Support","R&D","HR",
    "Legal"),
  pct = c(22, 28, 18, 12, 10, 6, 4)
) |>
  mutate(dept = fct_reorder(dept, pct))

ggplot(df, aes(x = dept, y = pct)) +
  geom_col(fill = "#1565C0",
           width = 0.55) +
  geom_text(aes(label = paste0(
    pct, "%")),
    hjust = -0.2, size = 3,
    fontface = "bold") +
  coord_flip() +
  theme_minimal() +
  theme(panel.grid.major.y =
    element_blank()) +
  labs(x = NULL,
    y = "Budget Share (%)",
    title = "Engineering Leads
             at 28%")
\end{lstlisting}
    \end{column}
    \begin{column}{0.45\textwidth}
      \includegraphics[width=\textwidth]{figures_m07/r_redesign}

      \vspace{4pt}
      {\small Key moves:
      \texttt{fct\_reorder()} sorts by value;
      \texttt{coord\_flip()} makes bars horizontal;
      \texttt{geom\_text()} adds direct labels.}
    \end{column}
  \end{columns}
\end{frame}

% ================================================================
\begin{frame}[fragile]{Redesign 1 in Python}
  \begin{columns}[T]
    \begin{column}{0.52\textwidth}
\begin{lstlisting}[style=pythonstyle]
import matplotlib.pyplot as plt
import numpy as np

depts = ["HR","Legal","R&D",
  "Support","Sales","Marketing",
  "Engineering"]
pcts = [6, 4, 10, 12, 18, 22, 28]

# Sort ascending for barh
idx = np.argsort(pcts)
s_d = [depts[i] for i in idx]
s_p = [pcts[i] for i in idx]

fig, ax = plt.subplots()
colors = ["#1565C0"] * 7
colors[-1] = "#E53935"  # accent
ax.barh(s_d, s_p,
    color=colors, height=0.55)
for i, v in enumerate(s_p):
    ax.text(v + 0.3, i, f"{v}%",
        va="center", fontsize=9,
        fontweight="bold")

# Tufte cleanup
for sp in ["top","right","left"]:
    ax.spines[sp].set_visible(False)
ax.set_xticks([])
ax.tick_params(left=False)
\end{lstlisting}
    \end{column}
    \begin{column}{0.45\textwidth}
      \includegraphics[width=\textwidth]{figures_m07/r_redesign}

      \vspace{4pt}
      {\small Python equivalent: sort with
      \texttt{np.argsort()}, use
      \texttt{ax.barh()}, remove spines
      explicitly.}
    \end{column}
  \end{columns}
\end{frame}

% ================================================================
\begin{frame}{Redesign 2: Dual Y-Axis $\rightarrow$ Indexed Lines}
  \begin{center}
    \includegraphics[width=0.92\textwidth]{figures_m07/redesign_dual_axis}
  \end{center}
  \begin{exampleblock}{Data Insight}
    \textbf{Diagnosis}: Integrity (implied false correlation). \\
    \textbf{Fix}: Index both series to 100 at month 1 so they share a
    common scale. Relative change becomes directly comparable.
  \end{exampleblock}
\end{frame}

% ================================================================
\begin{frame}[fragile]{Redesign 2 in R}
  \begin{columns}[T]
    \begin{column}{0.52\textwidth}
\begin{lstlisting}[style=rstyle]
# Index to 100 at first period
df_indexed <- df |>
  group_by(metric) |>
  mutate(
    index = value / first(value) * 100
  )

ggplot(df_indexed,
  aes(x = month, y = index,
      color = metric)) +
  geom_line(linewidth = 1.2) +
  geom_hline(yintercept = 100,
    linetype = "dashed",
    color = "grey60") +

  # Direct labels at endpoint
  geom_text(
    data = df_indexed |>
      filter(month == max(month)),
    aes(label = metric),
    nudge_x = 0.5, size = 3,
    fontface = "bold") +
  theme_minimal() +
  theme(legend.position = "none") +
  labs(y = "Index (M1 = 100)")
\end{lstlisting}
    \end{column}
    \begin{column}{0.45\textwidth}
      \includegraphics[width=\textwidth]{figures_m07/py_redesign}

      \vspace{4pt}
      {\small Indexing formula:
      $\text{index}_t = \frac{v_t}{v_1} \times 100$

      The reference line at 100 anchors
      ``no change'' visually.}
    \end{column}
  \end{columns}
\end{frame}

% ================================================================
\begin{frame}[fragile]{Redesign 2 in Python}
  \begin{columns}[T]
    \begin{column}{0.52\textwidth}
\begin{lstlisting}[style=pythonstyle]
# Index to 100 at month 1
rev_idx = revenue / revenue[0] * 100
usr_idx = users / users[0] * 100

fig, ax = plt.subplots()
ax.plot(months, rev_idx,
    color="#1565C0", lw=2)
ax.plot(months, usr_idx,
    color="#E53935", lw=2)

# Reference line
ax.axhline(y=100, color="#888",
    lw=0.8, linestyle="--")

# Direct labels
ax.text(12.2, rev_idx[-1],
    "Revenue", fontsize=8,
    fontweight="bold",
    color="#1565C0", va="center")
ax.text(12.2, usr_idx[-1],
    "Users", fontsize=8,
    fontweight="bold",
    color="#E53935", va="center")

ax.set_xlim(1, 14)
ax.set_ylabel("Index (M1=100)")
\end{lstlisting}
    \end{column}
    \begin{column}{0.45\textwidth}
      \includegraphics[width=\textwidth]{figures_m07/py_redesign}

      \vspace{4pt}
      {\small NumPy vectorised indexing
      replaces the \texttt{group\_by} +
      \texttt{mutate} pattern in R.}
    \end{column}
  \end{columns}
\end{frame}

% ================================================================
\begin{frame}{Redesign 3: Spaghetti Lines $\rightarrow$ Small Multiples}
  \begin{center}
    \includegraphics[width=0.92\textwidth]{figures_m07/redesign_spaghetti}
  \end{center}
  \begin{exampleblock}{Data Insight}
    \textbf{Diagnosis}: Encoding (6 overlapping lines → conjunction search). \\
    \textbf{Fix}: \texttt{facet\_wrap()} / \texttt{subplots(sharey=True)}
    separates series by proximity (Gestalt). Shared y-scale ensures
    fair cross-panel comparison.
  \end{exampleblock}
\end{frame}

% ================================================================
\begin{frame}{Redesign 4: Rainbow Scatter $\rightarrow$ Grey + Accent}
  \begin{center}
    \includegraphics[width=0.92\textwidth]{figures_m07/redesign_scatter}
  \end{center}
  \begin{exampleblock}{Data Insight}
    \textbf{Diagnosis}: Encoding (5 hues, no hierarchy) + Communication
    (no focal point). \\
    \textbf{Fix}: Grey all non-focal points, accent the target category,
    add a regression line for the focus group only.
  \end{exampleblock}
\end{frame}

% ================================================================
\begin{frame}{Redesign 5: Truncated Axis $\rightarrow$ Honest Baseline}
  \begin{center}
    \includegraphics[width=0.92\textwidth]{figures_m07/redesign_truncated}
  \end{center}
  \begin{alertblock}{Pro-Tip}
    \textbf{Diagnosis}: Integrity (lie factor $\approx$ 17). \\
    \textbf{Fix}: Start y-axis at zero, add a dashed reference line at the
    first-period value, and label each bar directly. The visual impression
    shifts from ``massive growth'' to ``steady improvement.''
  \end{alertblock}
\end{frame}

% ================================================================
\section{Peer Review Protocol}
% ================================================================

\begin{frame}{The Five-Step Peer Review}
  \begin{center}
    \includegraphics[width=0.78\textwidth]{figures_m07/peer_review}
  \end{center}
\end{frame}

% ================================================================
\begin{frame}{Giving Constructive Feedback}
  \begin{columns}[T]
    \begin{column}{0.48\textwidth}
      \textbf{Structure (SBI Model):}
      \begin{itemize}
        \item \textbf{S}ituation: ``In the Q3 report chart\ldots''
        \item \textbf{B}ehavior: ``\ldots the y-axis starts at 95,
              creating a lie factor of $\sim$17.''
        \item \textbf{I}mpact: ``This makes a 3\% change look like
              a 50\% change to leadership.''
      \end{itemize}
    \end{column}
    \begin{column}{0.48\textwidth}
      \textbf{Always include:}
      \begin{itemize}
        \item Which quadrant the issue falls in
        \item A specific, measurable diagnostic
              (lie factor, channel ranking)
        \item A concrete alternative (``try a
              zero-based axis with direct labels'')
      \end{itemize}
    \end{column}
  \end{columns}

  \vspace{6pt}
  \begin{alertblock}{Pro-Tip}
    Lead with what works well before identifying problems. Frame
    improvements as ``what if'' alternatives, not ``you should'' commands.
  \end{alertblock}
\end{frame}

% ================================================================
\begin{frame}{In-Class Exercise: Critique \& Redesign}
  \textbf{Instructions (20 minutes):}
  \begin{enumerate}
    \item Form pairs. Each person shows one chart from their own work
          or a published source.
    \item Apply the \textbf{five-step peer review} protocol.
    \item For each chart, identify:
          \begin{itemize}
            \item The \textbf{strongest} quadrant (what works)
            \item The \textbf{weakest} quadrant (what needs improvement)
          \end{itemize}
    \item Sketch a redesign on paper (3 minutes per chart).
    \item Share one redesign with the class (2 minutes per pair).
  \end{enumerate}

  \vspace{4pt}
  \begin{exampleblock}{Data Insight}
    The ability to critique constructively is the single most transferable
    skill from a visualization course. It applies to every chart you
    encounter for the rest of your career.
  \end{exampleblock}
\end{frame}

% ================================================================
\section{Synthesis}
% ================================================================

\begin{frame}{Five Redesign Patterns --- Summary}
  \centering
  \begin{tabular}{lll}
    \toprule
    \textbf{Problem} & \textbf{Quadrant} & \textbf{Fix} \\
    \midrule
    3-D pie chart       & Encoding + Craft & Horizontal bar \\
    Dual y-axis         & Integrity        & Indexed lines \\
    Spaghetti lines     & Encoding         & Small multiples \\
    Rainbow scatter     & Encoding + Comm. & Grey + accent \\
    Truncated axis      & Integrity        & Zero baseline \\
    \bottomrule
  \end{tabular}

  \vspace{8pt}
  Each pattern maps directly to a principle from Modules 1--6.
  The critique framework provides the diagnostic; the redesign
  implements the cure.
\end{frame}

% ================================================================
\begin{frame}{Key Takeaways}
  \begin{enumerate}
    \item The \textbf{four-quadrant framework} (integrity, encoding, craft,
          communication) provides a systematic, non-subjective critique
          language.
    \item Every redesign follows: \textbf{identify → diagnose → sketch →
          code → validate}.
    \item Five high-frequency problems have canonical fixes: pie → bar,
          dual axis → index, spaghetti → facets, rainbow → grey-accent,
          truncated → zero-based.
    \item \textbf{Peer review} using the SBI model (Situation, Behaviour,
          Impact) ensures feedback is specific and constructive.
  \end{enumerate}
\end{frame}

% ================================================================
\begin{frame}{Essential Readings for This Module}
  \begin{itemize}
    \item Knaflic, C.~N. (2015). \emph{Storytelling with Data},
          Chapter 3 (Clutter Is Your Enemy).
    \item Schwabish, J. (2021). \emph{Better Data Visualizations},
          Chapters 7--8 (Redesign Gallery).
    \item Wainer, H. (1984). ``How to Display Data Badly.''
          \emph{The American Statistician}, 38(2), 137--147.
    \item Kirk, A. (2019). \emph{Data Visualisation}, Chapter 9
          (Interactivity and Annotation).
  \end{itemize}

  \vspace{6pt}
  \textbf{Next Module}: R Environment Setup \& Base Graphics (W01-M08)
\end{frame}

\end{document}
